\section{3 结果}

\subsection{3.1 固-液-气三相碳污染物描述性统计}

本研究对校园污水管网中固相、液相和气相碳污染物进行了系统测定。气相碳污染物方面，CH₄浓度变化范围为0.63~2052.39 ppm（均值89.72±339.99 ppm，n=40），变异系数高达378.9\%，呈极强的空间异质性；CO₂浓度范围为566.08~19234.54 ppm（均值2003.55±3261.23 ppm，n=39），变异系数162.8\%，同样表现出高变异性；N₂O浓度相对较低且稳定，范围为0.38~1.24 ppm（均值0.45±0.19 ppm，n=20），变异系数41.3\%。VOCs浓度范围为238.00~869.00 ppb（均值512.62±160.54 ppb，n=40），属中等变异水平。

液相碳污染物中，TOC浓度范围为8.98~124.00 mg/L（均值44.52±30.95 mg/L，n=18），变异系数69.5\%；TC浓度范围为43.46~203.28 mg/L（均值114.87±44.58 mg/L，n=18），变异系数38.8\%；IC浓度范围为25.61~129.30 mg/L（均值70.34±27.69 mg/L，n=18），变异系数39.4\%。COD浓度变化范围最为悬殊，为14.06~12139.18 mg/L（均值1130.34±2802.51 mg/L，n=18），变异系数高达247.9\%，反映了不同功能区排水有机负荷的巨大差异。液相营养盐指标中，TN均值46.79±39.39 mg/L，NH₄⁺均值36.72±37.84 mg/L（变异系数103.1\%），NO₃⁻均值0.88±1.15 mg/L（变异系数130.4\%）。DO浓度范围为0.95~10.45 mg/L（均值4.80±2.26 mg/L，n=18），表明管网中存在明显的溶解氧梯度。pH范围为6.48~9.25（均值7.79±0.85），水温范围为8.10~20.30 ℃（均值15.42±3.30 ℃）。

固相碳污染物因采样量有限（n=2），仅作初步表征。固相总碳范围为95.27~299.46 g/kg（均值197.36 g/kg），有机碳范围为35.72~205.91 g/kg（均值120.81 g/kg），无机碳范围为8.41~9.15 g/kg（均值8.78 g/kg）。固相DOC变化范围为151.79~6513.93 mg/kg，变异系数135.0\%，提示沉积物中溶解态有机碳存在显著空间分异。

\subsection{3.2 正态性检验}

Shapiro-Wilk检验结果表明，25个变量中有11个符合正态分布（p>0.05），14个不符合正态分布（p≤0.05）。正态分布变量包括VOCs（p=0.352）、DO（p=0.360）、NaCl（p=0.078）、电导率（p=0.532）、水温（p=0.453）、pH（p=0.362）、TC（p=0.293）、IC（p=0.714）、TP（p=0.072）等。非正态分布变量主要包括CH₄、N₂O、CO₂、CO₂本底值、VOCs本底值、气温、TOC、TN、NH₄⁺、NO₃⁻、COD等。气相温室气体指标普遍呈非正态分布，与其高度偏态的频率分布特征一致。基于此，后续组间比较对非正态变量采用Mann-Whitney U非参数检验，对正态变量采用独立样本t检验。

\subsection{3.3 冬春季节组间差异}

组间比较结果显示，22个变量中有9个在冬春两季之间存在显著差异（表X）。气相碳污染物中，CH₄浓度春季（177.07±470.33 ppm）显著高于冬季（2.38±0.43 ppm）（Mann-Whitney U，p=0.019），差异近74倍，是所有变量中季节差异最为悬殊的指标。CO₂本底值冬季（638.00±21.71 ppm）显著高于春季（550.15±84.89 ppm）（p<0.001）。VOCs冬季（598.50±153.20 ppb）显著高于春季（426.75±117.94 ppb）（p<0.001），其本底值同样呈冬高春低趋势（p<0.001）。

液相指标中，COD浓度冬季（2426.38±3948.24 mg/L）显著高于春季（93.51±153.24 mg/L）（p<0.001），差异达一个数量级以上。NO₃⁻冬季（1.36±1.13 mg/L）显著高于春季（0.49±1.05 mg/L）（p=0.018）。温度指标方面，气温春季（16.05±1.23 ℃）显著高于冬季（6.65±6.42 ℃）（p<0.001），水温春季（17.24±1.66 ℃）显著高于冬季（13.15±3.51 ℃）（p=0.005）。

，液相TOC、TN、NH₄⁺、DO、pH等指标的冬春差异均未达到统计显著水平（p>0.05），表明这些水质指标在季节尺度上相对稳定。

\subsection{3.4 变量间相关性分析}

Pearson相关性分析表明多组显著的变量关联（图X）。气相指标内部，CO₂与NO₂呈极显著正相关（r=0.922，p<0.001），为所有相关系数中的最高值，提示两者可能受共同的环境因子驱动。N₂O与NaCl呈显著正相关（r=0.803，p<0.05），N₂O与TP亦呈显著正相关（r=0.716，p<0.05），暗示盐度和营养盐水平可能影响N₂O的产生过程。

液相碳氮指标之间，TOC与TC呈显著正相关（r=0.789，p<0.05），印证了有机碳是总碳的主要组分。TN与NH₄⁺呈极显著正相关（r=0.985，p<0.001），表明管网污水中氮素以铵态氮为主要存在形态，反映了厌氧环境下氨化作用的优势地位。

变量与采样时间的相关分析中，CO₂本底值与采样时间呈显著负相关（r=-0.622，p<0.001），VOCs本底值与采样时间呈极显著负相关（r=-0.893，p<0.001），而气温与采样时间呈显著正相关（r=0.712，p<0.001），共同反映了冬春季节更替的环境信号。电导率与COD呈显著负相关（r=-0.627，p<0.01），水温与IC呈显著负相关（r=-0.609，p<0.01）。

\subsection{3.5 碳平衡初步分析}

基于固-液-气三相碳浓度数据的初步碳平衡分析表明，液相碳（TC均值114.87 mg/L）是管网碳污染物的主要赋存形态，其中有机碳（TOC）占总碳的38.8\%，无机碳（IC）占61.2\%。固相沉积物中总碳含量较高（均值197.36 g/kg），有机碳占比约61.2\%（均值120.81 g/kg），无机碳占比约4.5\%（均值8.78 g/kg），表明沉积物是碳的重要储库。气相碳以CO₂为主要组分（均值2003.55 ppm），CH₄浓度虽然均值较低但具有极高的峰值（最大值2052.39 ppm），提示管网局部存在强烈的甲烷产生热点。三相碳的分配格局表明，液相和固相是碳的主要赋存相态，气相碳排放虽然占总碳比例相对较小，但其温室效应贡献不可忽视，尤其是CH₄的全球增温潜势为CO₂的28倍。